\section{Aplicación de los Algoritmos al Problema de Portfolio}
\label{sec:aplicacion}

En este apartado se recogen los elementos comunes a todos los algoritmos empleados en la práctica. La idea es fijar primero cómo se representa una solución, cómo se evalúa y cómo se mantiene la factibilidad; después se describen los operadores compartidos de selección, cruce y mutación.

\subsection{Esquema de representación}
Se utiliza una \textbf{representación real} basada en un vector de pesos
$$
S = (s_1, s_2, \dots, s_N)
$$
donde cada componente $s_i$ indica la fracción del capital invertida en el activo $i$. Una solución es factible si se cumplen simultáneamente estas condiciones:
\begin{itemize}
    \item $\sum_{i=1}^N s_i = 1$.
    \item Para cada activo, o bien $s_i = 0$, o bien $s_i \in [lo, hi]$.
\end{itemize}

La población inicial se genera aleatoriamente y, a continuación, se aplica reparación para garantizar que todos los individuos parten de soluciones válidas.

\subsection{Función objetivo}
La calidad de una cartera se mide con una función de \textit{fitness} que combina beneficio y riesgo. Como en el problema se busca maximizar, una mayor aptitud indica una mejor solución\footnote{El cambio que se ha introducido es que se ha realizado la raíz para el cálculo del riesgo.}.

\begin{algorithm}[H]
\SetAlgoLined
\KwIn{Solución $S$, beneficios logarítmicos $logBenefits$, matriz de covarianza $Cov$, parámetro $\lambda$}
\KwOut{Valor de \textit{fitness} de $S$}

$beneficio \leftarrow 0$

\For{$i \leftarrow 1$ \KwTo $N$}{
    $beneficio \leftarrow beneficio + s_i \cdot logBenefits_i$
}

$riesgo \leftarrow 0$

\For{$i \leftarrow 1$ \KwTo $N$}{
    \For{$j \leftarrow 1$ \KwTo $N$}{
        $riesgo \leftarrow riesgo + s_i \cdot s_j \cdot Cov(i,j)$
    }
}

\Return $beneficio - \lambda \cdot \sqrt{\max(0, riesgo)}$
\caption{Evaluación de una solución}
\end{algorithm}
\textit{Nota: Se incorpora el operador $\max(0, \cdot)$ dentro de la raíz cuadrada para garantizar la estabilidad numérica del algoritmo, evitando valores negativos infinitesimales producidos por errores de precisión en punto flotante (IEEE 754).}

\subsection{Operador de reparación}
Los operadores genéticos pueden producir soluciones no factibles, por lo que después de cada cruce o mutación se aplica un operador de reparación. Su misión es devolver el vector al espacio de búsqueda permitido sin alterar innecesariamente la estructura de la solución\footnote{Se usa el mismo pseudocódigo de la práctica 1 ya que no ha sufrido cambios.}.

\newpage
\begin{algorithm}[H]
\SetAlgoLined
\KwIn{Vector de pesos inválido $w$, límites $lo$ y $hi$, tolerancia $\epsilon$}
\KwOut{Vector de pesos $w$ factible}
 \tcp{Fase 1: Normalización inicial}
 $suma \leftarrow \sum_{i=1}^{N} w_i$\;
 \If{$suma > 0$ y $|suma - 1.0| > \epsilon$}{
    $w_i \leftarrow w_i / suma \quad \forall i$\;
 }
 \tcp{Fase 2: Poda estricta de límites}
 \For{$i \leftarrow 1$ \KwTo $N$}{
    \If{$w_i < lo$}{ $w_i \leftarrow 0.0$ }
    \ElseIf{$w_i > hi$}{ $w_i \leftarrow hi$ }
 }
 \tcp{Fase 3: Redistribución iterativa}
 $suma \leftarrow \sum_{i=1}^{N} w_i$\;
 \While{$|suma - 1.0| > \epsilon$}{
    $\delta \leftarrow 1.0 - suma$\;
    \eIf{$\delta > 0$}{
        \tcp{Falta capital: repartir entre activos sin saturar}
        $C \leftarrow \{i \mid 0 < w_i < hi\}$\;
        \eIf{$C \neq \emptyset$}{
            $incremento \leftarrow \delta / |C|$\;
            \For{$i \in C$}{
                $w_i \leftarrow w_i + \min(incremento, hi - w_i)$\;
            }
        }{
            \tcp{Retorno: no hay activos con margen}
            Seleccionar un índice $k$ aleatorio donde $w_k == 0$\;
            $w_k \leftarrow lo$\;
        }
    }{
        \tcp{Sobra capital: recortar de activos activos}
        $C \leftarrow \{i \mid w_i > lo\}$\;
        $decremento \leftarrow |\delta| / |C|$\;
        \For{$i \in C$}{
            $w_i \leftarrow w_i - \min(decremento, w_i - lo)$\;
        }
    }
    $suma \leftarrow \sum_{i=1}^{N} w_i$\;
 }
 \Return $w$\;
 \caption{Operador de Reparación \texttt{fix(w)}}
 \label{alg:fix}
\end{algorithm}


\subsection{Operadores comunes}
\subsubsection{Selección de padres: torneo de tamaño 3}
La selección se realiza mediante un torneo de tamaño $k=3$. Para elegir un padre se toman tres individuos al azar de la población y se conserva el que tenga mejor \textit{fitness}. Este operador introduce presión selectiva sin eliminar por completo la diversidad.

\begin{algorithm}[H]
\SetAlgoLined
\KwIn{Población $P$}
\KwOut{Índice del mejor individuo del torneo}

Elegir tres índices aleatorios de $P$ con reemplazo\;
Comparar sus valores de \textit{fitness}\;
Devolver el índice con mejor aptitud\;
\caption{Torneo de tamaño 3}
\end{algorithm}

\subsubsection{Operadores de cruce}
Se emplean dos operadores de cruce para representación real. En ambos casos se generan dos descendientes a partir de dos padres.

\begin{itemize}
    \item \textbf{Cruce aritmético aleatorio:} para cada gen $i$, se genera un coeficiente aleatorio $\sigma_i \in [0,1]$ y se construyen los hijos como combinaciones convexas de los padres.
    \item \textbf{Cruce BLX-$\alpha$:} para cada gen $i$, se toma el intervalo expandido definido por los valores de los padres y se muestrean los descendientes en ese rango. En esta práctica se usa $\alpha = 0{,}3$.
\end{itemize}

\begin{algorithm}[H]
\SetAlgoLined
\KwIn{Padres $P^1$ y $P^2$}
\KwOut{Hijos $H^1$ y $H^2$}

\For{cada posición $i$}{
    Generar $\sigma_i \in [0,1]$\;
    $H^1_i \leftarrow \sigma_i \cdot P^1_i + (1-\sigma_i) \cdot P^2_i$\;
    $H^2_i \leftarrow \sigma_i \cdot P^2_i + (1-\sigma_i) \cdot P^1_i$\;
}
\caption{Cruce aritmético}
\end{algorithm}

\begin{algorithm}[H]
\SetAlgoLined
\KwIn{Padres $P^1$ y $P^2$, parámetro $\alpha$}
\KwOut{Hijos $H^1$ y $H^2$}

\For{cada posición $i$}{
    $C_{min} \leftarrow \min(P^1_i, P^2_i)$\;
    $C_{max} \leftarrow \max(P^1_i, P^2_i)$\;
    $I \leftarrow C_{max} - C_{min}$\;
    Definir el intervalo $[C_{min} - \alpha I,\, C_{max} + \alpha I]$\;
    Generar $H^1_i$ y $H^2_i$ independientemente dentro de ese intervalo\;
}
\caption{Cruce BLX-$\alpha$}
\end{algorithm}

\subsubsection{Operador de mutación}
La mutación se implementa como un traslado parcial de capital entre dos genes distintos. Con probabilidad por individuo, se selecciona un origen y un destino y se transfiere un porcentaje fijo del valor del primero al segundo. En esta práctica el porcentaje transferido es del 15\%.

\begin{algorithm}[H]
\SetAlgoLined
\KwIn{Solución $S$}
\KwOut{Solución mutada y reparada}

Generar una probabilidad aleatoria para decidir si se muta el individuo\;
\If{no se activa la mutación}{
    \Return $S$
}
\If{la solución tiene menos de dos genes}{
    \Return $S$
}

Elegir dos posiciones distintas $i$ y $j$ al azar\;
$amount \leftarrow s_i \cdot 0.15$\;
$s_i \leftarrow s_i - amount$\;
$s_j \leftarrow s_j + amount$\;
Aplicar reparación sobre la solución\;
\caption{Mutación por transferencia}
\end{algorithm}